\begin{frame}{``정확도''만으로는 안 되는 이유: 정확도의 함정}
  \begin{itemize}
    \item 암 검진 모델: 전체 환자의 \textbf{99\%가 정상}라면,
      ``무조건 정상''이라고만 찍는 모델도 \textbf{정확도 99\%}
      --- 하지만 암 환자를 \textbf{단 한 명도} 잡아내지 못하는
      \textbf{쓸모없는 모델}.
  \end{itemize}
  \vskip0.6em
  \begin{center}
  \small
  \begin{tabular}{@{}lll@{}}
    \toprule
     & 환자 1000명 중 10명(1\%) 실제 암 & 모델: 무조건 ``정상'' \\
    \midrule
    정확도 & $990/1000 = \textbf{99\%}$ (매우 높아 보임) & \\
    Recall & $0/10 = \textbf{0\%}$ & 실제 암을 \textbf{단 한 명도}
      못 잡아 \\
    \bottomrule
  \end{tabular}
  \end{center}
  \vskip0.6em
  \kq{``정확도의 함정(accuracy paradox)'': 클래스 불균형 상황에서
  정확도라는 지표 하나만 봐도 훌륭해 보이지만 실제로는 쓸모가 없다.}
  $\to$ \kb{Precision / Recall / F1} 은 정확도가 \textbf{숨기는
  진실}을 드러내기 위한 지표.
\end{frame}
